% !TeX root = ../main.tex

\chapter{Evaluation}\label{chapter:evaluation}

% ~6--8 pages

\section{Experimental Setup}\label{sec:eval-setup}
\begin{itemize}
  \item Hardware: specify CPU, GPU (model, VRAM), RAM
  \item Software: Python version, PyTorch version, ONNX Runtime version, Docker version, Triton Inference Server version
  \item Primary PLM: ESM2-T6-8M (8 million parameters) --- chosen for fast iteration while retaining transformer behavior
  \item Baseline encodings: one-hot encoding, BLOSUM62
  \item Docker Compose configuration: server, 4 workers, Redis, PostgreSQL, SeaweedFS, Triton
  \item All experiments run on the same machine for fair comparison; report wall-clock times
\end{itemize}

\section{Test Coverage and Mutation Testing Results}\label{sec:eval-mutation}
\begin{itemize}
  \item Total test count: $\sim$180 test functions across 24 files
  \item Code coverage (line coverage, branch coverage) for unit tests --- report numbers with \texttt{pytest-cov}
  \item Mutation testing results:
    \begin{itemize}
      \item Total mutants generated by mutmut
      \item Killed vs.\ survived mutants
      \item Mutation score (\%)
      \item Which types of mutants survived --- what does this reveal about test suite weaknesses?
      \item Comparison: mutation score vs.\ line coverage --- are they correlated?
    \end{itemize}
  \item Table idea: mutation score breakdown by module (embeddings, predictions, projections, etc.)
\end{itemize}

\section{CI Execution Time Analysis}\label{sec:eval-time}
\begin{itemize}
  \item Wall-clock execution time per test category:
    \begin{itemize}
      \item Unit tests (with mocks): expected $\sim$seconds
      \item Integration tests (with FixedEmbedder): expected $\sim$minutes
      \item Integration tests (with real ESM2): expected $\sim$longer
      \item Property/oracle tests: time per oracle
      \item Performance tests: time per benchmark
      \item Metamorphic experiments (basic): time for 4-sequence variants
      \item Metamorphic experiments (UniRef50 scale): time for 1,750-sequence variants
    \end{itemize}
  \item \textbf{Key analysis}: plot fault-detection value (mutation score contribution, bug-finding capability) against CI time cost
  \item Trade-off curve: diminishing returns --- at what point does adding more property tests stop being worth the CI time?
  \item FixedEmbedder speedup: how much faster is CI with fixed vs.\ real embedder? What coverage is sacrificed?
  \item Recommendation: optimal test portfolio for different CI budgets (e.g., per-commit: unit + fixed integration; nightly: full suite; weekly: UniRef50 experiments)
  \item Figure idea: stacked bar chart of CI time per layer; scatter plot of time vs.\ coverage contribution
\end{itemize}

\section{Metamorphic Experiment Results}\label{sec:eval-metamorphic}

\subsection{X-Masking Results}\label{sec:eval-xmasking}
\begin{itemize}
  \item Cosine divergence profiles: plot cosine distance vs.\ masking ratio for ESM2
  \item Critical ratio $r^*$: at what masking percentage does cosine distance exceed 0.1?
  \item Monotonicity analysis: percentage of sequences with non-monotonic divergence; characterize violations
  \item ESM2 vs.\ one-hot comparison: one-hot should show linear divergence (bag-of-features); ESM2 shows non-linear (contextual model)
  \item ESM2 vs.\ random baseline: do random sequences diverge differently than real proteins?
  \item UniRef50 length-bin breakdown: does divergence profile depend on sequence length?
  \item Figure ideas: cosine similarity curve (ESM2 + one-hot + random), per-bin breakdown, L2 distance dual panel
\end{itemize}

\subsection{Reversed Sequence Results}\label{sec:eval-reversal}
\begin{itemize}
  \item Distribution of cosine distances for ESM2 reversed sequences --- confirm order-sensitivity
  \item One-hot baseline: confirm near-zero distance (order-insensitive, as expected)
  \item Double-reversal identity: confirm recovery within tolerance
  \item Interpretation: positional encoding in ESM2 captures sequence directionality
\end{itemize}

\subsection{AA Mutation Sensitivity Results}\label{sec:eval-mutation-sensitivity}
\begin{itemize}
  \item Heatmap: amino acid $\times$ masking ratio showing cosine distance
  \item Sensitivity ranking: which AAs cause largest/smallest embedding changes at $\sim$10\% mutation?
  \item Physicochemical grouping: charged (D, E, K, R) vs.\ hydrophobic (A, V, L, I) vs.\ polar (S, T, N, Q) vs.\ special (G, P, C, W)
  \item Comparison to X-masking: is replacing with a real AA more or less disruptive than replacing with X?
  \item UniRef50 scale: do sensitivity patterns hold across diverse real proteins?
  \item Figure ideas: heatmap, bar chart by physicochemical group, mutation vs.\ X-masking overlay
\end{itemize}

\section{Oracle Test Findings}\label{sec:eval-oracles}
\begin{itemize}
  \item Batch invariance: report cosine distances (expect $\ll 10^{-6}$); any violations?
  \item Prediction determinism: confirm bit-identical ONNX outputs across runs
  \item Output validity: any probabilities outside $[0,1]$ or not summing to 1? Any NaN/Inf?
  \item Shape invariance: any dimension mismatches?
  \item Projection distance preservation: Spearman correlation values for PCA, UMAP, t-SNE
  \item Projection determinism: PCA exact; UMAP/t-SNE Procrustes distances
  \item Table idea: oracle name, metric, threshold, observed value, pass/fail
\end{itemize}

\section{Performance Results}\label{sec:eval-performance}
\begin{itemize}
  \item Latency by sequence length: plot latency vs.\ sequence length (10, 15, 211, 400\,aa)
  \item Throughput: sequences per second for batch processing
  \item Memory: footprint per embedding, leak detection results
  \item O($n$) scaling: confirm linear scaling with sequence length (ratio $< 3\times$)
  \item Batch efficiency: batch vs.\ sequential processing speedup
  \item ESM2 vs.\ FixedEmbedder: orders-of-magnitude speed difference quantified
  \item Table/figure idea: latency table, memory bar chart, scaling line plot
\end{itemize}

\section{Discussion of Results}\label{sec:eval-discussion}
\begin{itemize}
  \item \textbf{Answering the RQ}: What is the trade-off between CI execution time and test coverage?
    \begin{itemize}
      \item Unit tests: fastest, catch validation/contract bugs, limited for ML-specific issues
      \item Integration tests (FixedEmbedder): moderate time, catch real flow bugs without GPU --- best bang-for-buck
      \item Oracle/property tests: moderate time, catch invariant violations that unit tests miss
      \item Metamorphic experiments: slow (especially UniRef50), provide unique behavioral insights but diminishing returns for CI
      \item Performance tests: moderate time, essential for regression detection
    \end{itemize}
  \item Recommended CI strategy: per-commit (unit + fixed integration), nightly (+ oracles + performance), weekly/release (+ metamorphic + real GPU)
  \item FixedEmbedder as enabler: makes the pyramid viable by decoupling CI from GPU availability
  \item Surprising findings (if any): unexpected oracle failures, non-monotonic masking behavior, etc.
\end{itemize}
